\makeatletter
\newif\ifBookIndexing
\BookIndexingfalse
\newcommand{\BookBoldPage}[1]{\textbf{\hyperpage{#1}}}
\DeclareRobustCommand{\BookRecordIndex}[1]{%
  \ifBookIndexing
    \expandafter\index\expandafter{\csname BookIndex@#1\endcsname|hyperpage}%
  \fi}
\DeclareRobustCommand{\BookIndexDefinition}[1]{%
  \ifBookIndexing
    \expandafter\index\expandafter{\csname BookIndex@#1\endcsname|BookBoldPage}%
  \fi}
\DeclareRobustCommand{\BookTerm}[2]{%
  \BookRecordIndex{#1}%
  \begingroup\let\BookRecordIndex\@gobble #2\endgroup}
\AddToHook{cmd/@outputpage/before}{\let\BookRecordIndex\@gobble}
\pdfstringdefDisableCommands{%
  \def\BookTerm#1#2{#2}%
  \let\BookRecordIndex\@gobble
  \let\BookIndexDefinition\@gobble}
\expandafter\def\csname BookIndex@T0001\endcsname{xinxingshi@辛形式（symplectic form）}
\expandafter\let\csname BookOriginal@termSymplecticForm\endcsname\termSymplecticForm
\renewcommand{\termSymplecticForm}{\BookTerm{T0001}{\csname BookOriginal@termSymplecticForm\endcsname}}
\expandafter\def\csname BookIndex@T0002\endcsname{hamidunweifentongpei@哈密顿微分同胚（Hamiltonian diffeomorphism）}
\expandafter\let\csname BookOriginal@termHamiltonianDiffeomorphism\endcsname\termHamiltonianDiffeomorphism
\renewcommand{\termHamiltonianDiffeomorphism}{\BookTerm{T0002}{\csname BookOriginal@termHamiltonianDiffeomorphism\endcsname}}
\expandafter\def\csname BookIndex@T0003\endcsname{lagelangriziliuxing@拉格朗日子流形（Lagrangian submanifold）}
\expandafter\let\csname BookOriginal@termLagrangianSubmanifold\endcsname\termLagrangianSubmanifold
\renewcommand{\termLagrangianSubmanifold}{\BookTerm{T0003}{\csname BookOriginal@termLagrangianSubmanifold\endcsname}}
\expandafter\def\csname BookIndex@T0004\endcsname{weiquanchunquxian@伪全纯曲线（pseudoholomorphic curve）}
\expandafter\let\csname BookOriginal@termPseudoholomorphicCurve\endcsname\termPseudoholomorphicCurve
\renewcommand{\termPseudoholomorphicCurve}{\BookTerm{T0004}{\csname BookOriginal@termPseudoholomorphicCurve\endcsname}}
\expandafter\def\csname BookIndex@T0005\endcsname{daikaileiliuxing@殆凯勒流形（almost Kähler manifold）}
\expandafter\let\csname BookOriginal@termAlmostKahlerManifold\endcsname\termAlmostKahlerManifold
\renewcommand{\termAlmostKahlerManifold}{\BookTerm{T0005}{\csname BookOriginal@termAlmostKahlerManifold\endcsname}}
\expandafter\def\csname BookIndex@T0006\endcsname{huofeierfanshu@霍费尔范数（Hofer norm）}
\expandafter\let\csname BookOriginal@termHoferNorm\endcsname\termHoferNorm
\renewcommand{\termHoferNorm}{\BookTerm{T0006}{\csname BookOriginal@termHoferNorm\endcsname}}
\expandafter\def\csname BookIndex@T0007\endcsname{lagelangrijiaofuluoertongdiao@拉格朗日交弗洛尔同调（Lagrangian intersection Floer homology）}
\expandafter\let\csname BookOriginal@termLagrangianIntersectionFloerHomology\endcsname\termLagrangianIntersectionFloerHomology
\renewcommand{\termLagrangianIntersectionFloerHomology}{\BookTerm{T0007}{\csname BookOriginal@termLagrangianIntersectionFloerHomology\endcsname}}
\expandafter\def\csname BookIndex@T0008\endcsname{lagelangrijiaolilun@拉格朗日交理论（Lagrangian intersection theory）}
\expandafter\let\csname BookOriginal@termLagrangianIntersectionTheory\endcsname\termLagrangianIntersectionTheory
\renewcommand{\termLagrangianIntersectionTheory}{\BookTerm{T0008}{\csname BookOriginal@termLagrangianIntersectionTheory\endcsname}}
\expandafter\def\csname BookIndex@T0009\endcsname{fuluoerfuxing@弗洛尔复形（Floer complex）}
\expandafter\let\csname BookOriginal@termFloerComplex\endcsname\termFloerComplex
\renewcommand{\termFloerComplex}{\BookTerm{T0009}{\csname BookOriginal@termFloerComplex\endcsname}}
\expandafter\def\csname BookIndex@T0010\endcsname{luguo@滤过（filtration）}
\expandafter\let\csname BookOriginal@termFiltration\endcsname\termFiltration
\renewcommand{\termFiltration}{\BookTerm{T0010}{\csname BookOriginal@termFiltration\endcsname}}
\expandafter\def\csname BookIndex@T0011\endcsname{zuoyongliangjifen@作用量积分（action integral）}
\expandafter\let\csname BookOriginal@termActionIntegral\endcsname\termActionIntegral
\renewcommand{\termActionIntegral}{\BookTerm{T0011}{\csname BookOriginal@termActionIntegral\endcsname}}
\expandafter\def\csname BookIndex@T0012\endcsname{maopaoxianxiang@冒泡现象（bubbling phenomenon）}
\expandafter\let\csname BookOriginal@termBubblingPhenomenon\endcsname\termBubblingPhenomenon
\renewcommand{\termBubblingPhenomenon}{\BookTerm{T0012}{\csname BookOriginal@termBubblingPhenomenon\endcsname}}
\expandafter\def\csname BookIndex@T0013\endcsname{lingjiemian@零截面（zero section）}
\expandafter\let\csname BookOriginal@termZeroSection\endcsname\termZeroSection
\renewcommand{\termZeroSection}{\BookTerm{T0013}{\csname BookOriginal@termZeroSection\endcsname}}
\expandafter\def\csname BookIndex@T0014\endcsname{hamidunguidao@哈密顿轨道（Hamiltonian trajectory）}
\expandafter\let\csname BookOriginal@termHamiltonianTrajectory\endcsname\termHamiltonianTrajectory
\renewcommand{\termHamiltonianTrajectory}{\BookTerm{T0014}{\csname BookOriginal@termHamiltonianTrajectory\endcsname}}
\expandafter\def\csname BookIndex@T0015\endcsname{jibenshengchenghanshu@基本生成函数（basic generating function）}
\expandafter\let\csname BookOriginal@termBasicGeneratingFunction\endcsname\termBasicGeneratingFunction
\renewcommand{\termBasicGeneratingFunction}{\BookTerm{T0015}{\csname BookOriginal@termBasicGeneratingFunction\endcsname}}
\expandafter\def\csname BookIndex@T0016\endcsname{ziranbianjietiaojian@自然边界条件（natural boundary condition）}
\expandafter\let\csname BookOriginal@termNaturalBoundaryCondition\endcsname\termNaturalBoundaryCondition
\renewcommand{\termNaturalBoundaryCondition}{\BookTerm{T0016}{\csname BookOriginal@termNaturalBoundaryCondition\endcsname}}
\expandafter\def\csname BookIndex@T0017\endcsname{xianxinghuasuanzi@线性化算子（linearization operator）}
\expandafter\let\csname BookOriginal@termLinearizationOperator\endcsname\termLinearizationOperator
\renewcommand{\termLinearizationOperator}{\BookTerm{T0017}{\csname BookOriginal@termLinearizationOperator\endcsname}}
\expandafter\def\csname BookIndex@T0018\endcsname{tuiqian@推前（push forward）}
\expandafter\let\csname BookOriginal@termPushForward\endcsname\termPushForward
\renewcommand{\termPushForward}{\BookTerm{T0018}{\csname BookOriginal@termPushForward\endcsname}}
\expandafter\def\csname BookIndex@T0019\endcsname{youxianweiyuehua@有限维约化（finite dimensional reduction）}
\expandafter\let\csname BookOriginal@termFiniteDimensionalReduction\endcsname\termFiniteDimensionalReduction
\renewcommand{\termFiniteDimensionalReduction}{\BookTerm{T0019}{\csname BookOriginal@termFiniteDimensionalReduction\endcsname}}
\expandafter\def\csname BookIndex@T0020\endcsname{mokongjian@模空间（moduli space）}
\expandafter\let\csname BookOriginal@termModuliSpace\endcsname\termModuliSpace
\renewcommand{\termModuliSpace}{\BookTerm{T0020}{\csname BookOriginal@termModuliSpace\endcsname}}
\expandafter\def\csname BookIndex@T0021\endcsname{kexilimanxing@柯西--黎曼型（Cauchy–Riemann type）}
\expandafter\let\csname BookOriginal@termCauchyRiemannType\endcsname\termCauchyRiemannType
\renewcommand{\termCauchyRiemannType}{\BookTerm{T0021}{\csname BookOriginal@termCauchyRiemannType\endcsname}}
\expandafter\def\csname BookIndex@T0022\endcsname{fanziduiouyangmiersifangcheng@反自对偶杨--米尔斯方程（anti self dual yang mills equation）}
\expandafter\let\csname BookOriginal@termAntiSelfDualYangMillsEquation\endcsname\termAntiSelfDualYangMillsEquation
\renewcommand{\termAntiSelfDualYangMillsEquation}{\BookTerm{T0022}{\csname BookOriginal@termAntiSelfDualYangMillsEquation\endcsname}}
\expandafter\def\csname BookIndex@T0023\endcsname{saibogeweitengdanjifangcheng@塞伯格--威腾单极方程（seiberg Witten monopole equation）}
\expandafter\let\csname BookOriginal@termSeibergWittenMonopoleEquation\endcsname\termSeibergWittenMonopoleEquation
\renewcommand{\termSeibergWittenMonopoleEquation}{\BookTerm{T0023}{\csname BookOriginal@termSeibergWittenMonopoleEquation\endcsname}}
\expandafter\def\csname BookIndex@T0024\endcsname{hengjiexing@横截性（transversality）}
\expandafter\let\csname BookOriginal@termTransversality\endcsname\termTransversality
\renewcommand{\termTransversality}{\BookTerm{T0024}{\csname BookOriginal@termTransversality\endcsname}}
\expandafter\def\csname BookIndex@T0025\endcsname{jihejiaodianshu@几何交点数（geometric intersection number）}
\expandafter\let\csname BookOriginal@termGeometricIntersectionNumber\endcsname\termGeometricIntersectionNumber
\renewcommand{\termGeometricIntersectionNumber}{\BookTerm{T0025}{\csname BookOriginal@termGeometricIntersectionNumber\endcsname}}
\expandafter\def\csname BookIndex@T0026\endcsname{dianfandaifujiegou@典范殆复结构（canonical almost complex structure）}
\expandafter\let\csname BookOriginal@termCanonicalAlmostComplexStructure\endcsname\termCanonicalAlmostComplexStructure
\renewcommand{\termCanonicalAlmostComplexStructure}{\BookTerm{T0026}{\csname BookOriginal@termCanonicalAlmostComplexStructure\endcsname}}
\expandafter\def\csname BookIndex@T0027\endcsname{danshebanjing@单射半径（injectivity radius）}
\expandafter\let\csname BookOriginal@termInjectivityRadius\endcsname\termInjectivityRadius
\renewcommand{\termInjectivityRadius}{\BookTerm{T0027}{\csname BookOriginal@termInjectivityRadius\endcsname}}
\expandafter\def\csname BookIndex@T0028\endcsname{fanxinfanshe@反辛反射（anti symplectic reflection）}
\expandafter\let\csname BookOriginal@termAntiSymplecticReflection\endcsname\termAntiSymplecticReflection
\renewcommand{\termAntiSymplecticReflection}{\BookTerm{T0028}{\csname BookOriginal@termAntiSymplecticReflection\endcsname}}
\expandafter\def\csname BookIndex@T0029\endcsname{zhijixiantapu@直极限拓扑（direct limit topology）}
\expandafter\let\csname BookOriginal@termDirectLimitTopology\endcsname\termDirectLimitTopology
\renewcommand{\termDirectLimitTopology}{\BookTerm{T0029}{\csname BookOriginal@termDirectLimitTopology\endcsname}}
\expandafter\def\csname BookIndex@T0030\endcsname{raodongkexilimanfangcheng@扰动柯西--黎曼方程（perturbed Cauchy–Riemann equation）}
\expandafter\let\csname BookOriginal@termPerturbedCauchyRiemannEquation\endcsname\termPerturbedCauchyRiemannEquation
\renewcommand{\termPerturbedCauchyRiemannEquation}{\BookTerm{T0030}{\csname BookOriginal@termPerturbedCauchyRiemannEquation\endcsname}}
\expandafter\def\csname BookIndex@T0031\endcsname{youxiannengliang@有限能量（finite energy）}
\expandafter\let\csname BookOriginal@termFiniteEnergy\endcsname\termFiniteEnergy
\renewcommand{\termFiniteEnergy}{\BookTerm{T0031}{\csname BookOriginal@termFiniteEnergy\endcsname}}
\expandafter\def\csname BookIndex@T0032\endcsname{youjiexiang@有界像（bounded image）}
\expandafter\let\csname BookOriginal@termBoundedImage\endcsname\termBoundedImage
\renewcommand{\termBoundedImage}{\BookTerm{T0032}{\csname BookOriginal@termBoundedImage\endcsname}}
\expandafter\def\csname BookIndex@T0033\endcsname{xianyannengliangjie@先验能量界（apriori energy bound）}
\expandafter\let\csname BookOriginal@termAprioriEnergyBound\endcsname\termAprioriEnergyBound
\renewcommand{\termAprioriEnergyBound}{\BookTerm{T0033}{\csname BookOriginal@termAprioriEnergyBound\endcsname}}
\expandafter\def\csname BookIndex@T0034\endcsname{c0jie@$C^0$ 界（czero bound）}
\expandafter\let\csname BookOriginal@termCzeroBound\endcsname\termCzeroBound
\renewcommand{\termCzeroBound}{\BookTerm{T0034}{\csname BookOriginal@termCzeroBound\endcsname}}
\expandafter\def\csname BookIndex@T0035\endcsname{wuqiongyuanchudexunfuxing@无穷远处的驯服性（tameness at \ensuremath{\infty}）}
\expandafter\let\csname BookOriginal@termTamenessAtInfinity\endcsname\termTamenessAtInfinity
\renewcommand{\termTamenessAtInfinity}{\BookTerm{T0035}{\csname BookOriginal@termTamenessAtInfinity\endcsname}}
\expandafter\def\csname BookIndex@T0036\endcsname{xianweishuchengyingshe@纤维数乘映射（fiber multiplication map）}
\expandafter\let\csname BookOriginal@termFiberMultiplicationMap\endcsname\termFiberMultiplicationMap
\renewcommand{\termFiberMultiplicationMap}{\BookTerm{T0036}{\csname BookOriginal@termFiberMultiplicationMap\endcsname}}
\expandafter\def\csname BookIndex@T0037\endcsname{zijulunzheng@自举论证（bootstrap argument）}
\expandafter\let\csname BookOriginal@termBootstrapArgument\endcsname\termBootstrapArgument
\renewcommand{\termBootstrapArgument}{\BookTerm{T0037}{\csname BookOriginal@termBootstrapArgument\endcsname}}
\expandafter\def\csname BookIndex@T0038\endcsname{maopaofenxi@冒泡分析（bubbling off analysis）}
\expandafter\let\csname BookOriginal@termBubblingOffAnalysis\endcsname\termBubblingOffAnalysis
\renewcommand{\termBubblingOffAnalysis}{\BookTerm{T0038}{\csname BookOriginal@termBubblingOffAnalysis\endcsname}}
\expandafter\def\csname BookIndex@T0039\endcsname{fuleidehuomuxingzhi@弗雷德霍姆性质（Fredholm property）}
\expandafter\let\csname BookOriginal@termFredholmProperty\endcsname\termFredholmProperty
\renewcommand{\termFredholmProperty}{\BookTerm{T0039}{\csname BookOriginal@termFredholmProperty\endcsname}}
\expandafter\def\csname BookIndex@T0040\endcsname{fuluoerbianjieyingshe@弗洛尔边界映射（Floer boundary map）}
\expandafter\let\csname BookOriginal@termFloerBoundaryMap\endcsname\termFloerBoundaryMap
\renewcommand{\termFloerBoundaryMap}{\BookTerm{T0040}{\csname BookOriginal@termFloerBoundaryMap\endcsname}}
\expandafter\def\csname BookIndex@T0041\endcsname{lianyingshe@链映射（chain map）}
\expandafter\let\csname BookOriginal@termChainMap\endcsname\termChainMap
\renewcommand{\termChainMap}{\BookTerm{T0041}{\csname BookOriginal@termChainMap\endcsname}}
\expandafter\def\csname BookIndex@T0042\endcsname{jiechuxingbianjie@接触型边界（contact type boundary）}
\expandafter\let\csname BookOriginal@termContactTypeBoundary\endcsname\termContactTypeBoundary
\renewcommand{\termContactTypeBoundary}{\BookTerm{T0042}{\csname BookOriginal@termContactTypeBoundary\endcsname}}
\expandafter\def\csname BookIndex@T0043\endcsname{duozhongcidiaohehanshu@多重次调和函数（plurisubharmonic function）}
\expandafter\let\csname BookOriginal@termPlurisubharmonicFunction\endcsname\termPlurisubharmonicFunction
\renewcommand{\termPlurisubharmonicFunction}{\BookTerm{T0043}{\csname BookOriginal@termPlurisubharmonicFunction\endcsname}}
\expandafter\def\csname BookIndex@T0044\endcsname{qiangzuidazhiyuanli@强最大值原理（strong maximum principle）}
\expandafter\let\csname BookOriginal@termStrongMaximumPrinciple\endcsname\termStrongMaximumPrinciple
\renewcommand{\termStrongMaximumPrinciple}{\BookTerm{T0044}{\csname BookOriginal@termStrongMaximumPrinciple\endcsname}}
\expandafter\def\csname BookIndex@T0045\endcsname{zhishushuaijianguji@指数衰减估计（exponential decay estimate）}
\expandafter\let\csname BookOriginal@termExponentialDecayEstimate\endcsname\termExponentialDecayEstimate
\renewcommand{\termExponentialDecayEstimate}{\BookTerm{T0045}{\csname BookOriginal@termExponentialDecayEstimate\endcsname}}
\expandafter\def\csname BookIndex@T0046\endcsname{yuanpancong@圆盘丛（disc bundle）}
\expandafter\let\csname BookOriginal@termDiscBundle\endcsname\termDiscBundle
\renewcommand{\termDiscBundle}{\BookTerm{T0046}{\csname BookOriginal@termDiscBundle\endcsname}}
\expandafter\def\csname BookIndex@T0047\endcsname{leirangdeziliuxing@勒让德子流形（legendrian submanifold）}
\expandafter\let\csname BookOriginal@termLegendrianSubmanifold\endcsname\termLegendrianSubmanifold
\renewcommand{\termLegendrianSubmanifold}{\BookTerm{T0047}{\csname BookOriginal@termLegendrianSubmanifold\endcsname}}
\expandafter\def\csname BookIndex@T0048\endcsname{jiechufenbu@接触分布（contact distribution）}
\expandafter\let\csname BookOriginal@termContactDistribution\endcsname\termContactDistribution
\renewcommand{\termContactDistribution}{\BookTerm{T0048}{\csname BookOriginal@termContactDistribution\endcsname}}
\expandafter\def\csname BookIndex@T0049\endcsname{jingxiangxiangliangchang@径向向量场（radial vector field）}
\expandafter\let\csname BookOriginal@termRadialVectorField\endcsname\termRadialVectorField
\renewcommand{\termRadialVectorField}{\BookTerm{T0049}{\csname BookOriginal@termRadialVectorField\endcsname}}
\expandafter\def\csname BookIndex@T0050\endcsname{tongyoucanshu@通有参数（generic parameter）}
\expandafter\let\csname BookOriginal@termGenericParameter\endcsname\termGenericParameter
\renewcommand{\termGenericParameter}{\BookTerm{T0050}{\csname BookOriginal@termGenericParameter\endcsname}}
\expandafter\def\csname BookIndex@T0051\endcsname{fuluoersuanzi@弗洛尔算子（Floer operator）}
\expandafter\let\csname BookOriginal@termFloerOperator\endcsname\termFloerOperator
\renewcommand{\termFloerOperator}{\BookTerm{T0051}{\csname BookOriginal@termFloerOperator\endcsname}}
\expandafter\def\csname BookIndex@T0052\endcsname{tonghenlei@同痕类（isotopy class）}
\expandafter\let\csname BookOriginal@termIsotopyClass\endcsname\termIsotopyClass
\renewcommand{\termIsotopyClass}{\BookTerm{T0052}{\csname BookOriginal@termIsotopyClass\endcsname}}
\expandafter\def\csname BookIndex@T0053\endcsname{fuleidehuomuzhengzexing@弗雷德霍姆正则性（Fredholm regularity）}
\expandafter\let\csname BookOriginal@termFredholmRegularity\endcsname\termFredholmRegularity
\renewcommand{\termFredholmRegularity}{\BookTerm{T0053}{\csname BookOriginal@termFredholmRegularity\endcsname}}
\expandafter\def\csname BookIndex@T0054\endcsname{dianfanfenci@典范分次（canonical grading）}
\expandafter\let\csname BookOriginal@termCanonicalGrading\endcsname\termCanonicalGrading
\renewcommand{\termCanonicalGrading}{\BookTerm{T0054}{\csname BookOriginal@termCanonicalGrading\endcsname}}
\expandafter\def\csname BookIndex@T0055\endcsname{xinglieshicong@行列式丛（determinant bundle）}
\expandafter\let\csname BookOriginal@termDeterminantBundle\endcsname\termDeterminantBundle
\renewcommand{\termDeterminantBundle}{\BookTerm{T0055}{\csname BookOriginal@termDeterminantBundle\endcsname}}
\expandafter\def\csname BookIndex@T0056\endcsname{fencifuxing@分次复形（graded complex）}
\expandafter\let\csname BookOriginal@termGradedComplex\endcsname\termGradedComplex
\renewcommand{\termGradedComplex}{\BookTerm{T0056}{\csname BookOriginal@termGradedComplex\endcsname}}
\expandafter\def\csname BookIndex@T0057\endcsname{yidongbianjietiaojian@移动边界条件（moving boundary condition）}
\expandafter\let\csname BookOriginal@termMovingBoundaryCondition\endcsname\termMovingBoundaryCondition
\renewcommand{\termMovingBoundaryCondition}{\BookTerm{T0057}{\csname BookOriginal@termMovingBoundaryCondition\endcsname}}
\expandafter\def\csname BookIndex@T0058\endcsname{canshuhuamokongjian@参数化模空间（parameterized moduli space）}
\expandafter\let\csname BookOriginal@termParameterizedModuliSpace\endcsname\termParameterizedModuliSpace
\renewcommand{\termParameterizedModuliSpace}{\BookTerm{T0058}{\csname BookOriginal@termParameterizedModuliSpace\endcsname}}
\expandafter\def\csname BookIndex@T0059\endcsname{liantonglunyingshe@链同伦映射（chain homotopy map）}
\expandafter\let\csname BookOriginal@termChainHomotopyMap\endcsname\termChainHomotopyMap
\renewcommand{\termChainHomotopyMap}{\BookTerm{T0059}{\csname BookOriginal@termChainHomotopyMap\endcsname}}
\expandafter\def\csname BookIndex@T0060\endcsname{zhanhedingli@粘合定理（gluing theorem）}
\expandafter\let\csname BookOriginal@termGluingTheorem\endcsname\termGluingTheorem
\renewcommand{\termGluingTheorem}{\BookTerm{T0060}{\csname BookOriginal@termGluingTheorem\endcsname}}
\expandafter\def\csname BookIndex@T0061\endcsname{moersimokongjian@莫尔斯模空间（Morse moduli space）}
\expandafter\let\csname BookOriginal@termMorseModuliSpace\endcsname\termMorseModuliSpace
\renewcommand{\termMorseModuliSpace}{\BookTerm{T0061}{\csname BookOriginal@termMorseModuliSpace\endcsname}}
\expandafter\def\csname BookIndex@T0062\endcsname{qiyishangtongdiao@奇异上同调（singular cohomology）}
\expandafter\let\csname BookOriginal@termSingularCohomology\endcsname\termSingularCohomology
\renewcommand{\termSingularCohomology}{\BookTerm{T0062}{\csname BookOriginal@termSingularCohomology\endcsname}}
\expandafter\def\csname BookIndex@T0063\endcsname{lagelangripububianliang@拉格朗日谱不变量（Lagrangian spectral invariant）}
\expandafter\let\csname BookOriginal@termLagrangianSpectralInvariant\endcsname\termLagrangianSpectralInvariant
\renewcommand{\termLagrangianSpectralInvariant}{\BookTerm{T0063}{\csname BookOriginal@termLagrangianSpectralInvariant\endcsname}}
\expandafter\def\csname BookIndex@T0064\endcsname{jixiaojidalilun@极小极大理论（minimax theory）}
\expandafter\let\csname BookOriginal@termMinimaxTheory\endcsname\termMinimaxTheory
\renewcommand{\termMinimaxTheory}{\BookTerm{T0064}{\csname BookOriginal@termMinimaxTheory\endcsname}}
\expandafter\def\csname BookIndex@T0065\endcsname{lianjiexingzhi@链接性质（linking property）}
\expandafter\let\csname BookOriginal@termLinkingProperty\endcsname\termLinkingProperty
\renewcommand{\termLinkingProperty}{\BookTerm{T0065}{\csname BookOriginal@termLinkingProperty\endcsname}}
\expandafter\def\csname BookIndex@T0066\endcsname{palaisimeiertiaojian@帕莱--斯梅尔条件（palais Smale condition）}
\expandafter\let\csname BookOriginal@termPalaisSmaleCondition\endcsname\termPalaisSmaleCondition
\renewcommand{\termPalaisSmaleCondition}{\BookTerm{T0066}{\csname BookOriginal@termPalaisSmaleCondition\endcsname}}
\expandafter\def\csname BookIndex@T0067\endcsname{xiangxialuguo@向下滤过（downward filtration）}
\expandafter\let\csname BookOriginal@termDownwardFiltration\endcsname\termDownwardFiltration
\renewcommand{\termDownwardFiltration}{\BookTerm{T0067}{\csname BookOriginal@termDownwardFiltration\endcsname}}
\expandafter\def\csname BookIndex@T0068\endcsname{xiangduitongdiaoqun@相对同调群（relative homology group）}
\expandafter\let\csname BookOriginal@termRelativeHomologyGroup\endcsname\termRelativeHomologyGroup
\renewcommand{\termRelativeHomologyGroup}{\BookTerm{T0068}{\csname BookOriginal@termRelativeHomologyGroup\endcsname}}
\expandafter\def\csname BookIndex@T0069\endcsname{zuoyongliangcengji@作用量层级（action level）}
\expandafter\let\csname BookOriginal@termActionLevel\endcsname\termActionLevel
\renewcommand{\termActionLevel}{\BookTerm{T0069}{\csname BookOriginal@termActionLevel\endcsname}}
\expandafter\def\csname BookIndex@T0070\endcsname{dengdulianxu@等度连续（equicontinuous）}
\expandafter\let\csname BookOriginal@termEquicontinuous\endcsname\termEquicontinuous
\renewcommand{\termEquicontinuous}{\BookTerm{T0070}{\csname BookOriginal@termEquicontinuous\endcsname}}
\expandafter\def\csname BookIndex@T0071\endcsname{fuluoerfangchengdexingbian@弗洛尔方程的形变（deformation of Floer equations）}
\expandafter\let\csname BookOriginal@termDeformationOfFloerEquations\endcsname\termDeformationOfFloerEquations
\renewcommand{\termDeformationOfFloerEquations}{\BookTerm{T0071}{\csname BookOriginal@termDeformationOfFloerEquations\endcsname}}
\expandafter\def\csname BookIndex@T0072\endcsname{donglixuebanben@动力学版本（dynamical version）}
\expandafter\let\csname BookOriginal@termDynamicalVersion\endcsname\termDynamicalVersion
\renewcommand{\termDynamicalVersion}{\BookTerm{T0072}{\csname BookOriginal@termDynamicalVersion\endcsname}}
\expandafter\def\csname BookIndex@T0073\endcsname{jihebanben@几何版本（geometric version）}
\expandafter\let\csname BookOriginal@termGeometricVersion\endcsname\termGeometricVersion
\renewcommand{\termGeometricVersion}{\BookTerm{T0073}{\csname BookOriginal@termGeometricVersion\endcsname}}
\expandafter\def\csname BookIndex@T0074\endcsname{baochiluguodetonggou@保持滤过的同构（filtration preserving isomorphism）}
\expandafter\let\csname BookOriginal@termFiltrationPreservingIsomorphism\endcsname\termFiltrationPreservingIsomorphism
\renewcommand{\termFiltrationPreservingIsomorphism}{\BookTerm{T0074}{\csname BookOriginal@termFiltrationPreservingIsomorphism\endcsname}}
\expandafter\def\csname BookIndex@T0075\endcsname{youxiaozuoyongliangfanhan@有效作用量泛函（effective action functional）}
\expandafter\let\csname BookOriginal@termEffectiveActionFunctional\endcsname\termEffectiveActionFunctional
\renewcommand{\termEffectiveActionFunctional}{\BookTerm{T0075}{\csname BookOriginal@termEffectiveActionFunctional\endcsname}}
\expandafter\def\csname BookIndex@T0076\endcsname{bianjiegongxian@边界贡献（boundary contribution）}
\expandafter\let\csname BookOriginal@termBoundaryContribution\endcsname\termBoundaryContribution
\renewcommand{\termBoundaryContribution}{\BookTerm{T0076}{\csname BookOriginal@termBoundaryContribution\endcsname}}
\expandafter\def\csname BookIndex@T0077\endcsname{guanzhuanglinyu@管状邻域（tubular neighborhood）}
\expandafter\let\csname BookOriginal@termTubularNeighborhood\endcsname\termTubularNeighborhood
\renewcommand{\termTubularNeighborhood}{\BookTerm{T0077}{\csname BookOriginal@termTubularNeighborhood\endcsname}}
\expandafter\def\csname BookIndex@T0078\endcsname{dabuwensitanlinyu@达布--温斯坦邻域（Darboux weinstein neighborhood）}
\expandafter\let\csname BookOriginal@termDarbouxWeinsteinNeighborhood\endcsname\termDarbouxWeinsteinNeighborhood
\renewcommand{\termDarbouxWeinsteinNeighborhood}{\BookTerm{T0078}{\csname BookOriginal@termDarbouxWeinsteinNeighborhood\endcsname}}
\expandafter\def\csname BookIndex@T0079\endcsname{hamidunpububianliang@哈密顿谱不变量（Hamiltonian spectral invariant）}
\expandafter\let\csname BookOriginal@termHamiltonianSpectralInvariant\endcsname\termHamiltonianSpectralInvariant
\renewcommand{\termHamiltonianSpectralInvariant}{\BookTerm{T0079}{\csname BookOriginal@termHamiltonianSpectralInvariant\endcsname}}
\expandafter\def\csname BookIndex@T0080\endcsname{jieduanhanshu@截断函数（cut off function）}
\expandafter\let\csname BookOriginal@termCutOffFunction\endcsname\termCutOffFunction
\renewcommand{\termCutOffFunction}{\BookTerm{T0080}{\csname BookOriginal@termCutOffFunction\endcsname}}
\expandafter\def\csname BookIndex@T0081\endcsname{xiangweihanshu@相位函数（phase function）}
\expandafter\let\csname BookOriginal@termPhaseFunction\endcsname\termPhaseFunction
\renewcommand{\termPhaseFunction}{\BookTerm{T0081}{\csname BookOriginal@termPhaseFunction\endcsname}}
\expandafter\def\csname BookIndex@T0082\endcsname{weijubufenxi@微局部分析（microlocal analysis）}
\expandafter\let\csname BookOriginal@termMicrolocalAnalysis\endcsname\termMicrolocalAnalysis
\renewcommand{\termMicrolocalAnalysis}{\BookTerm{T0082}{\csname BookOriginal@termMicrolocalAnalysis\endcsname}}
\expandafter\def\csname BookIndex@T0083\endcsname{fuliyejifensuanzi@傅里叶积分算子（fourier integral operator）}
\expandafter\let\csname BookOriginal@termFourierIntegralOperator\endcsname\termFourierIntegralOperator
\renewcommand{\termFourierIntegralOperator}{\BookTerm{T0083}{\csname BookOriginal@termFourierIntegralOperator\endcsname}}
\expandafter\def\csname BookIndex@T0084\endcsname{tapuhamidunliang@拓扑哈密顿量（topological Hamiltonian）}
\expandafter\let\csname BookOriginal@termTopologicalHamiltonian\endcsname\termTopologicalHamiltonian
\renewcommand{\termTopologicalHamiltonian}{\BookTerm{T0084}{\csname BookOriginal@termTopologicalHamiltonian\endcsname}}
\expandafter\def\csname BookIndex@T0085\endcsname{hamiduntapu@哈密顿拓扑（Hamiltonian topology）}
\expandafter\let\csname BookOriginal@termHamiltonianTopology\endcsname\termHamiltonianTopology
\renewcommand{\termHamiltonianTopology}{\BookTerm{T0085}{\csname BookOriginal@termHamiltonianTopology\endcsname}}
\expandafter\def\csname BookIndex@T0086\endcsname{bandanbufen@半单部分（semi simple part）}
\expandafter\let\csname BookOriginal@termSemiSimplePart\endcsname\termSemiSimplePart
\renewcommand{\termSemiSimplePart}{\BookTerm{T0086}{\csname BookOriginal@termSemiSimplePart\endcsname}}
\expandafter\def\csname BookIndex@T0087\endcsname{lipuxicitapu@利普希茨拓扑（lipschitz topology）}
\expandafter\let\csname BookOriginal@termLipschitzTopology\endcsname\termLipschitzTopology
\renewcommand{\termLipschitzTopology}{\BookTerm{T0087}{\csname BookOriginal@termLipschitzTopology\endcsname}}
\expandafter\def\csname BookIndex@T0088\endcsname{fencengziliuxing@分层子流形（stratified submanifold）}
\expandafter\let\csname BookOriginal@termStratifiedSubmanifold\endcsname\termStratifiedSubmanifold
\renewcommand{\termStratifiedSubmanifold}{\BookTerm{T0088}{\csname BookOriginal@termStratifiedSubmanifold\endcsname}}
\expandafter\def\csname BookIndex@T0089\endcsname{huofeierjuli@霍费尔距离（Hofer distance）}
\expandafter\let\csname BookOriginal@termHoferDistance\endcsname\termHoferDistance
\renewcommand{\termHoferDistance}{\BookTerm{T0089}{\csname BookOriginal@termHoferDistance\endcsname}}
\expandafter\def\csname BookIndex@T0090\endcsname{hamidunxingbian@哈密顿形变（Hamiltonian deformation）}
\expandafter\let\csname BookOriginal@termHamiltonianDeformation\endcsname\termHamiltonianDeformation
\renewcommand{\termHamiltonianDeformation}{\BookTerm{T0090}{\csname BookOriginal@termHamiltonianDeformation\endcsname}}
\expandafter\def\csname BookIndex@T0091\endcsname{duojiemianlagelangriziliuxing@多截面拉格朗日子流形（multisection Lagrangian）}
\expandafter\let\csname BookOriginal@termMultisectionLagrangian\endcsname\termMultisectionLagrangian
\renewcommand{\termMultisectionLagrangian}{\BookTerm{T0091}{\csname BookOriginal@termMultisectionLagrangian\endcsname}}
\expandafter\def\csname BookIndex@T0092\endcsname{jiandian@尖点（cusp point）}
\expandafter\let\csname BookOriginal@termCuspPoint\endcsname\termCuspPoint
\renewcommand{\termCuspPoint}{\BookTerm{T0092}{\csname BookOriginal@termCuspPoint\endcsname}}
\expandafter\def\csname BookIndex@T0093\endcsname{limanluohegongshi@黎曼--罗赫公式（Riemann roch formula）}
\expandafter\let\csname BookOriginal@termRiemannRochFormula\endcsname\termRiemannRochFormula
\renewcommand{\termRiemannRochFormula}{\BookTerm{T0093}{\csname BookOriginal@termRiemannRochFormula\endcsname}}
\expandafter\def\csname BookIndex@T0094\endcsname{xuweishu@虚维数（virtual dimension）}
\expandafter\let\csname BookOriginal@termVirtualDimension\endcsname\termVirtualDimension
\renewcommand{\termVirtualDimension}{\BookTerm{T0094}{\csname BookOriginal@termVirtualDimension\endcsname}}
\expandafter\def\csname BookIndex@T0095\endcsname{quanchunyingshe@全纯映射（holomorphic map）}
\expandafter\let\csname BookOriginal@termHolomorphicMap\endcsname\termHolomorphicMap
\renewcommand{\termHolomorphicMap}{\BookTerm{T0095}{\csname BookOriginal@termHolomorphicMap\endcsname}}
\expandafter\def\csname BookIndex@T0096\endcsname{diyichenshu@第一陈数（first Chern number）}
\expandafter\let\csname BookOriginal@termFirstChernNumber\endcsname\termFirstChernNumber
\renewcommand{\termFirstChernNumber}{\BookTerm{T0096}{\csname BookOriginal@termFirstChernNumber\endcsname}}
\expandafter\def\csname BookIndex@T0097\endcsname{masiluofulei@马斯洛夫类（Maslov class）}
\expandafter\let\csname BookOriginal@termMaslovClass\endcsname\termMaslovClass
\renewcommand{\termMaslovClass}{\BookTerm{T0097}{\csname BookOriginal@termMaslovClass\endcsname}}
\expandafter\def\csname BookIndex@T0098\endcsname{lagelangrizikongjian@拉格朗日子空间（Lagrangian subspace）}
\expandafter\let\csname BookOriginal@termLagrangianSubspace\endcsname\termLagrangianSubspace
\renewcommand{\termLagrangianSubspace}{\BookTerm{T0098}{\csname BookOriginal@termLagrangianSubspace\endcsname}}
\expandafter\def\csname BookIndex@T0099\endcsname{qixingkongjian@齐性空间（homogeneous space）}
\expandafter\let\csname BookOriginal@termHomogeneousSpace\endcsname\termHomogeneousSpace
\renewcommand{\termHomogeneousSpace}{\BookTerm{T0099}{\csname BookOriginal@termHomogeneousSpace\endcsname}}
\expandafter\def\csname BookIndex@T0100\endcsname{fuchengfa@复乘法（complex multiplication）}
\expandafter\let\csname BookOriginal@termComplexMultiplication\endcsname\termComplexMultiplication
\renewcommand{\termComplexMultiplication}{\BookTerm{T0100}{\csname BookOriginal@termComplexMultiplication\endcsname}}
\expandafter\def\csname BookIndex@T0101\endcsname{tonglundengjia@同伦等价（homotopy equivalence）}
\expandafter\let\csname BookOriginal@termHomotopyEquivalence\endcsname\termHomotopyEquivalence
\renewcommand{\termHomotopyEquivalence}{\BookTerm{T0101}{\csname BookOriginal@termHomotopyEquivalence\endcsname}}
\expandafter\def\csname BookIndex@T0102\endcsname{shixinglei@示性类（characteristic class）}
\expandafter\let\csname BookOriginal@termCharacteristicClass\endcsname\termCharacteristicClass
\renewcommand{\termCharacteristicClass}{\BookTerm{T0102}{\csname BookOriginal@termCharacteristicClass\endcsname}}
\expandafter\def\csname BookIndex@T0103\endcsname{youjuzhen@酉矩阵（unitary matrix）}
\expandafter\let\csname BookOriginal@termUnitaryMatrix\endcsname\termUnitaryMatrix
\renewcommand{\termUnitaryMatrix}{\BookTerm{T0103}{\csname BookOriginal@termUnitaryMatrix\endcsname}}
\expandafter\def\csname BookIndex@T0104\endcsname{duichengyoujuzhen@对称酉矩阵（symmetric unitary matrix）}
\expandafter\let\csname BookOriginal@termSymmetricUnitaryMatrix\endcsname\termSymmetricUnitaryMatrix
\renewcommand{\termSymmetricUnitaryMatrix}{\BookTerm{T0104}{\csname BookOriginal@termSymmetricUnitaryMatrix\endcsname}}
\expandafter\def\csname BookIndex@T0105\endcsname{huanraoshu@环绕数（winding number）}
\expandafter\let\csname BookOriginal@termWindingNumber\endcsname\termWindingNumber
\renewcommand{\termWindingNumber}{\BookTerm{T0105}{\csname BookOriginal@termWindingNumber\endcsname}}
\expandafter\def\csname BookIndex@T0106\endcsname{guangyimasiluofuzhibiao@广义马斯洛夫指标（generalized Maslov index）}
\expandafter\let\csname BookOriginal@termGeneralizedMaslovIndex\endcsname\termGeneralizedMaslovIndex
\renewcommand{\termGeneralizedMaslovIndex}{\BookTerm{T0106}{\csname BookOriginal@termGeneralizedMaslovIndex\endcsname}}
\expandafter\def\csname BookIndex@T0107\endcsname{quanshizicong@全实子丛（totally real subbundle）}
\expandafter\let\csname BookOriginal@termTotallyRealSubbundle\endcsname\termTotallyRealSubbundle
\renewcommand{\termTotallyRealSubbundle}{\BookTerm{T0107}{\csname BookOriginal@termTotallyRealSubbundle\endcsname}}
\expandafter\def\csname BookIndex@T0108\endcsname{xinxiangliangcong@辛向量丛（symplectic vector bundle）}
\expandafter\let\csname BookOriginal@termSymplecticVectorBundle\endcsname\termSymplecticVectorBundle
\renewcommand{\termSymplecticVectorBundle}{\BookTerm{T0108}{\csname BookOriginal@termSymplecticVectorBundle\endcsname}}
\expandafter\def\csname BookIndex@T0109\endcsname{lagelangrizicong@拉格朗日子丛（Lagrangian subbundle）}
\expandafter\let\csname BookOriginal@termLagrangianSubbundle\endcsname\termLagrangianSubbundle
\renewcommand{\termLagrangianSubbundle}{\BookTerm{T0109}{\csname BookOriginal@termLagrangianSubbundle\endcsname}}
\expandafter\def\csname BookIndex@T0110\endcsname{pingfanhua@平凡化（trivialization）}
\expandafter\let\csname BookOriginal@termTrivialization\endcsname\termTrivialization
\renewcommand{\termTrivialization}{\BookTerm{T0110}{\csname BookOriginal@termTrivialization\endcsname}}
\expandafter\def\csname BookIndex@T0111\endcsname{xinpingfanhua@辛平凡化（symplectic trivialization）}
\expandafter\let\csname BookOriginal@termSymplecticTrivialization\endcsname\termSymplecticTrivialization
\renewcommand{\termSymplecticTrivialization}{\BookTerm{T0111}{\csname BookOriginal@termSymplecticTrivialization\endcsname}}
\expandafter\def\csname BookIndex@T0112\endcsname{fuhua@复化（complexification）}
\expandafter\let\csname BookOriginal@termComplexification\endcsname\termComplexification
\renewcommand{\termComplexification}{\BookTerm{T0112}{\csname BookOriginal@termComplexification\endcsname}}
\expandafter\def\csname BookIndex@T0113\endcsname{tonglunbubianxing@同伦不变性（homotopy invariance）}
\expandafter\let\csname BookOriginal@termHomotopyInvariance\endcsname\termHomotopyInvariance
\renewcommand{\termHomotopyInvariance}{\BookTerm{T0113}{\csname BookOriginal@termHomotopyInvariance\endcsname}}
\expandafter\def\csname BookIndex@T0114\endcsname{peibianbubianxing@配边不变性（cobordism invariance）}
\expandafter\let\csname BookOriginal@termCobordismInvariance\endcsname\termCobordismInvariance
\renewcommand{\termCobordismInvariance}{\BookTerm{T0114}{\csname BookOriginal@termCobordismInvariance\endcsname}}
\expandafter\def\csname BookIndex@T0115\endcsname{dingxiangpeibian@定向配边（oriented cobordism）}
\expandafter\let\csname BookOriginal@termOrientedCobordism\endcsname\termOrientedCobordism
\renewcommand{\termOrientedCobordism}{\BookTerm{T0115}{\csname BookOriginal@termOrientedCobordism\endcsname}}
\expandafter\def\csname BookIndex@T0116\endcsname{duobianxingyingshe@多边形映射（polygonal map）}
\expandafter\let\csname BookOriginal@termPolygonalMap\endcsname\termPolygonalMap
\renewcommand{\termPolygonalMap}{\BookTerm{T0116}{\csname BookOriginal@termPolygonalMap\endcsname}}
\expandafter\def\csname BookIndex@T0117\endcsname{tapuzhibiao@拓扑指标（topological index）}
\expandafter\let\csname BookOriginal@termTopologicalIndex\endcsname\termTopologicalIndex
\renewcommand{\termTopologicalIndex}{\BookTerm{T0117}{\csname BookOriginal@termTopologicalIndex\endcsname}}
\expandafter\def\csname BookIndex@T0118\endcsname{hengjielagelangrilian@横截拉格朗日链（transversal Lagrangian chain）}
\expandafter\let\csname BookOriginal@termTransversalLagrangianChain\endcsname\termTransversalLagrangianChain
\renewcommand{\termTransversalLagrangianChain}{\BookTerm{T0118}{\csname BookOriginal@termTransversalLagrangianChain\endcsname}}
\expandafter\def\csname BookIndex@T0119\endcsname{biaojidian@标记点（marked point）}
\expandafter\let\csname BookOriginal@termMarkedPoint\endcsname\termMarkedPoint
\renewcommand{\termMarkedPoint}{\BookTerm{T0119}{\csname BookOriginal@termMarkedPoint\endcsname}}
\expandafter\def\csname BookIndex@T0120\endcsname{chuankongyuanpan@穿孔圆盘（punctured disc）}
\expandafter\let\csname BookOriginal@termPuncturedDisc\endcsname\termPuncturedDisc
\renewcommand{\termPuncturedDisc}{\BookTerm{T0120}{\csname BookOriginal@termPuncturedDisc\endcsname}}
\expandafter\def\csname BookIndex@T0121\endcsname{gongxingtonggou@共形同构（conformal isomorphism）}
\expandafter\let\csname BookOriginal@termConformalIsomorphism\endcsname\termConformalIsomorphism
\renewcommand{\termConformalIsomorphism}{\BookTerm{T0121}{\csname BookOriginal@termConformalIsomorphism\endcsname}}
\expandafter\def\csname BookIndex@T0122\endcsname{hengjielian@横截链（transversal chain）}
\expandafter\let\csname BookOriginal@termTransversalChain\endcsname\termTransversalChain
\renewcommand{\termTransversalChain}{\BookTerm{T0122}{\csname BookOriginal@termTransversalChain\endcsname}}
\expandafter\def\csname BookIndex@T0123\endcsname{fenpianguanghualujing@分片光滑路径（piecewise smooth path）}
\expandafter\let\csname BookOriginal@termPiecewiseSmoothPath\endcsname\termPiecewiseSmoothPath
\renewcommand{\termPiecewiseSmoothPath}{\BookTerm{T0123}{\csname BookOriginal@termPiecewiseSmoothPath\endcsname}}
\expandafter\def\csname BookIndex@T0124\endcsname{daibianlimanqumian@带边黎曼曲面（bordered Riemann surface）}
\expandafter\let\csname BookOriginal@termBorderedRiemannSurface\endcsname\termBorderedRiemannSurface
\renewcommand{\termBorderedRiemannSurface}{\BookTerm{T0124}{\csname BookOriginal@termBorderedRiemannSurface\endcsname}}
\expandafter\def\csname BookIndex@T0125\endcsname{dabuzuobiaoka@达布坐标卡（Darboux chart）}
\expandafter\let\csname BookOriginal@termDarbouxChart\endcsname\termDarbouxChart
\renewcommand{\termDarbouxChart}{\BookTerm{T0125}{\csname BookOriginal@termDarbouxChart\endcsname}}
\expandafter\def\csname BookIndex@T0126\endcsname{gaosiyingshe@高斯映射（gauss map）}
\expandafter\let\csname BookOriginal@termGaussMap\endcsname\termGaussMap
\renewcommand{\termGaussMap}{\BookTerm{T0126}{\csname BookOriginal@termGaussMap\endcsname}}
\expandafter\def\csname BookIndex@T0127\endcsname{shengufanchou@深谷范畴（Fukaya category）}
\expandafter\let\csname BookOriginal@termFukayaCategory\endcsname\termFukayaCategory
\renewcommand{\termFukayaCategory}{\BookTerm{T0127}{\csname BookOriginal@termFukayaCategory\endcsname}}
\expandafter\def\csname BookIndex@T0128\endcsname{fuleidehuomuzhibiao@弗雷德霍姆指标（Fredholm index）}
\expandafter\let\csname BookOriginal@termFredholmIndex\endcsname\termFredholmIndex
\renewcommand{\termFredholmIndex}{\BookTerm{T0128}{\csname BookOriginal@termFredholmIndex\endcsname}}
\expandafter\def\csname BookIndex@T0129\endcsname{jizhundaolu@基准道路（base path）}
\expandafter\let\csname BookOriginal@termBasePath\endcsname\termBasePath
\renewcommand{\termBasePath}{\BookTerm{T0129}{\csname BookOriginal@termBasePath\endcsname}}
\expandafter\def\csname BookIndex@T0130\endcsname{wanyoufugaikongjian@万有覆盖空间（universal covering space）}
\expandafter\let\csname BookOriginal@termUniversalCoveringSpace\endcsname\termUniversalCoveringSpace
\renewcommand{\termUniversalCoveringSpace}{\BookTerm{T0130}{\csname BookOriginal@termUniversalCoveringSpace\endcsname}}
\expandafter\def\csname BookIndex@T0131\endcsname{huandaiyingshe@环带映射（annular map）}
\expandafter\let\csname BookOriginal@termAnnularMap\endcsname\termAnnularMap
\renewcommand{\termAnnularMap}{\BookTerm{T0131}{\csname BookOriginal@termAnnularMap\endcsname}}
\expandafter\def\csname BookIndex@T0132\endcsname{xinmianji@辛面积（symplectic area）}
\expandafter\let\csname BookOriginal@termSymplecticArea\endcsname\termSymplecticArea
\renewcommand{\termSymplecticArea}{\BookTerm{T0132}{\csname BookOriginal@termSymplecticArea\endcsname}}
\expandafter\def\csname BookIndex@T0133\endcsname{fugaibianhuan@覆盖变换（deck transformation）}
\expandafter\let\csname BookOriginal@termDeckTransformation\endcsname\termDeckTransformation
\renewcommand{\termDeckTransformation}{\BookTerm{T0133}{\csname BookOriginal@termDeckTransformation\endcsname}}
\expandafter\def\csname BookIndex@T0134\endcsname{qunhuan@群环（group ring）}
\expandafter\let\csname BookOriginal@termGroupRing\endcsname\termGroupRing
\renewcommand{\termGroupRing}{\BookTerm{T0134}{\csname BookOriginal@termGroupRing\endcsname}}
\expandafter\def\csname BookIndex@T0135\endcsname{xingshihe@形式和（formal sum）}
\expandafter\let\csname BookOriginal@termFormalSum\endcsname\termFormalSum
\renewcommand{\termFormalSum}{\BookTerm{T0135}{\csname BookOriginal@termFormalSum\endcsname}}
\expandafter\def\csname BookIndex@T0136\endcsname{fencirmo@分次 $R$-模（graded r module）}
\expandafter\let\csname BookOriginal@termGradedRModule\endcsname\termGradedRModule
\renewcommand{\termGradedRModule}{\BookTerm{T0136}{\csname BookOriginal@termGradedRModule\endcsname}}
\expandafter\def\csname BookIndex@T0137\endcsname{fugaitouying@覆盖投影（covering projection）}
\expandafter\let\csname BookOriginal@termCoveringProjection\endcsname\termCoveringProjection
\renewcommand{\termCoveringProjection}{\BookTerm{T0137}{\csname BookOriginal@termCoveringProjection\endcsname}}
\expandafter\def\csname BookIndex@T0138\endcsname{jubutisheng@局部提升（local lifting）}
\expandafter\let\csname BookOriginal@termLocalLifting\endcsname\termLocalLifting
\renewcommand{\termLocalLifting}{\BookTerm{T0138}{\csname BookOriginal@termLocalLifting\endcsname}}
\expandafter\def\csname BookIndex@T0139\endcsname{linjiedianji@临界点集（critical point set）}
\expandafter\let\csname BookOriginal@termCriticalPointSet\endcsname\termCriticalPointSet
\renewcommand{\termCriticalPointSet}{\BookTerm{T0139}{\csname BookOriginal@termCriticalPointSet\endcsname}}
\expandafter\def\csname BookIndex@T0140\endcsname{zhucong@主丛（principal bundle）}
\expandafter\let\csname BookOriginal@termPrincipalBundle\endcsname\termPrincipalBundle
\renewcommand{\termPrincipalBundle}{\BookTerm{T0140}{\csname BookOriginal@termPrincipalBundle\endcsname}}
\expandafter\def\csname BookIndex@T0141\endcsname{naozi@挠子（torsor）}
\expandafter\let\csname BookOriginal@termTorsor\endcsname\termTorsor
\renewcommand{\termTorsor}{\BookTerm{T0141}{\csname BookOriginal@termTorsor\endcsname}}
\expandafter\def\csname BookIndex@T0142\endcsname{jueduimoersizhibiao@绝对莫尔斯指标（absolute Morse index）}
\expandafter\let\csname BookOriginal@termAbsoluteMorseIndex\endcsname\termAbsoluteMorseIndex
\renewcommand{\termAbsoluteMorseIndex}{\BookTerm{T0142}{\csname BookOriginal@termAbsoluteMorseIndex\endcsname}}
\expandafter\def\csname BookIndex@T0143\endcsname{daijidiandelagelangriziliuxing@带基点的拉格朗日子流形（pointed Lagrangian submanifold）}
\expandafter\let\csname BookOriginal@termPointedLagrangianSubmanifold\endcsname\termPointedLagrangianSubmanifold
\renewcommand{\termPointedLagrangianSubmanifold}{\BookTerm{T0143}{\csname BookOriginal@termPointedLagrangianSubmanifold\endcsname}}
\expandafter\def\csname BookIndex@T0144\endcsname{mao@锚（anchor）}
\expandafter\let\csname BookOriginal@termAnchor\endcsname\termAnchor
\renewcommand{\termAnchor}{\BookTerm{T0144}{\csname BookOriginal@termAnchor\endcsname}}
\expandafter\def\csname BookIndex@T0145\endcsname{maoxitong@锚系统（anchor system）}
\expandafter\let\csname BookOriginal@termAnchorSystem\endcsname\termAnchorSystem
\renewcommand{\termAnchorSystem}{\BookTerm{T0145}{\csname BookOriginal@termAnchorSystem\endcsname}}
\expandafter\def\csname BookIndex@T0146\endcsname{zhongdiexingzhi@重叠性质（overlapping property）}
\expandafter\let\csname BookOriginal@termOverlappingProperty\endcsname\termOverlappingProperty
\renewcommand{\termOverlappingProperty}{\BookTerm{T0146}{\csname BookOriginal@termOverlappingProperty\endcsname}}
\expandafter\def\csname BookIndex@T0147\endcsname{dingjiedai@定界带（bounding band）}
\expandafter\let\csname BookOriginal@termBoundingBand\endcsname\termBoundingBand
\renewcommand{\termBoundingBand}{\BookTerm{T0147}{\csname BookOriginal@termBoundingBand\endcsname}}
\expandafter\def\csname BookIndex@T0148\endcsname{chouxiangfuluoerfuxing@抽象弗洛尔复形（abstract Floer complex）}
\expandafter\let\csname BookOriginal@termAbstractFloerComplex\endcsname\termAbstractFloerComplex
\renewcommand{\termAbstractFloerComplex}{\BookTerm{T0148}{\csname BookOriginal@termAbstractFloerComplex\endcsname}}
\expandafter\def\csname BookIndex@T0149\endcsname{fenciluguofuluoernuoweikefufuxing@分次滤过弗洛尔--诺维科夫复形（graded filtered Floer Novikov complex）}
\expandafter\let\csname BookOriginal@termGradedFilteredFloerNovikovComplex\endcsname\termGradedFilteredFloerNovikovComplex
\renewcommand{\termGradedFilteredFloerNovikovComplex}{\BookTerm{T0149}{\csname BookOriginal@termGradedFilteredFloerNovikovComplex\endcsname}}
\expandafter\def\csname BookIndex@T0150\endcsname{zhouqitongtai@周期同态（period homomorphism）}
\expandafter\let\csname BookOriginal@termPeriodHomomorphism\endcsname\termPeriodHomomorphism
\renewcommand{\termPeriodHomomorphism}{\BookTerm{T0150}{\csname BookOriginal@termPeriodHomomorphism\endcsname}}
\expandafter\def\csname BookIndex@T0151\endcsname{cishutongtai@次数同态（degree homomorphism）}
\expandafter\let\csname BookOriginal@termDegreeHomomorphism\endcsname\termDegreeHomomorphism
\renewcommand{\termDegreeHomomorphism}{\BookTerm{T0151}{\csname BookOriginal@termDegreeHomomorphism\endcsname}}
\expandafter\def\csname BookIndex@T0152\endcsname{fuluoerlianqun@弗洛尔链群（Floer chain group）}
\expandafter\let\csname BookOriginal@termFloerChainGroup\endcsname\termFloerChainGroup
\renewcommand{\termFloerChainGroup}{\BookTerm{T0152}{\csname BookOriginal@termFloerChainGroup\endcsname}}
\expandafter\def\csname BookIndex@T0153\endcsname{nuoweikefuyouxianxingtiaojian@诺维科夫有限性条件（Novikov finiteness condition）}
\expandafter\let\csname BookOriginal@termNovikovFinitenessCondition\endcsname\termNovikovFinitenessCondition
\renewcommand{\termNovikovFinitenessCondition}{\BookTerm{T0153}{\csname BookOriginal@termNovikovFinitenessCondition\endcsname}}
\expandafter\def\csname BookIndex@T0154\endcsname{cengjihanshu@层级函数（level function）}
\expandafter\let\csname BookOriginal@termLevelFunction\endcsname\termLevelFunction
\renewcommand{\termLevelFunction}{\BookTerm{T0154}{\csname BookOriginal@termLevelFunction\endcsname}}
\expandafter\def\csname BookIndex@T0155\endcsname{feiajimideduliang@非阿基米德度量（non archimedean metric）}
\expandafter\let\csname BookOriginal@termNonArchimedeanMetric\endcsname\termNonArchimedeanMetric
\renewcommand{\termNonArchimedeanMetric}{\BookTerm{T0155}{\csname BookOriginal@termNonArchimedeanMetric\endcsname}}
\expandafter\def\csname BookIndex@T0156\endcsname{luguolianfuxing@滤过链复形（filtered chain complex）}
\expandafter\let\csname BookOriginal@termFilteredChainComplex\endcsname\termFilteredChainComplex
\renewcommand{\termFilteredChainComplex}{\BookTerm{T0156}{\csname BookOriginal@termFilteredChainComplex\endcsname}}
\expandafter\def\csname BookIndex@T0157\endcsname{fuluoerlianmo@弗洛尔链模（Floer chain module）}
\expandafter\let\csname BookOriginal@termFloerChainModule\endcsname\termFloerChainModule
\renewcommand{\termFloerChainModule}{\BookTerm{T0157}{\csname BookOriginal@termFloerChainModule\endcsname}}
\expandafter\def\csname BookIndex@T0158\endcsname{feiajimidefanshu@非阿基米德范数（non archimedean norm）}
\expandafter\let\csname BookOriginal@termNonArchimedeanNorm\endcsname\termNonArchimedeanNorm
\renewcommand{\termNonArchimedeanNorm}{\BookTerm{T0158}{\csname BookOriginal@termNonArchimedeanNorm\endcsname}}
\expandafter\def\csname BookIndex@T0159\endcsname{l2duliang@$L^2$-度量（ltwo metric）}
\expandafter\let\csname BookOriginal@termLtwoMetric\endcsname\termLtwoMetric
\renewcommand{\termLtwoMetric}{\BookTerm{T0159}{\csname BookOriginal@termLtwoMetric\endcsname}}
\expandafter\def\csname BookIndex@T0160\endcsname{youjietiduliuxian@有界梯度流线（bounded gradient trajectory）}
\expandafter\let\csname BookOriginal@termBoundedGradientTrajectory\endcsname\termBoundedGradientTrajectory
\renewcommand{\termBoundedGradientTrajectory}{\BookTerm{T0160}{\csname BookOriginal@termBoundedGradientTrajectory\endcsname}}
\expandafter\def\csname BookIndex@T0161\endcsname{diaohenengliang@调和能量（harmonic energy）}
\expandafter\let\csname BookOriginal@termHarmonicEnergy\endcsname\termHarmonicEnergy
\renewcommand{\termHarmonicEnergy}{\BookTerm{T0161}{\csname BookOriginal@termHarmonicEnergy\endcsname}}
\expandafter\def\csname BookIndex@T0162\endcsname{fuluoermokongjian@弗洛尔模空间（Floer moduli space）}
\expandafter\let\csname BookOriginal@termFloerModuliSpace\endcsname\termFloerModuliSpace
\renewcommand{\termFloerModuliSpace}{\BookTerm{T0162}{\csname BookOriginal@termFloerModuliSpace\endcsname}}
\expandafter\def\csname BookIndex@T0163\endcsname{zhudianfanshu@逐点范数（pointwise norm）}
\expandafter\let\csname BookOriginal@termPointwiseNorm\endcsname\termPointwiseNorm
\renewcommand{\termPointwiseNorm}{\BookTerm{T0163}{\csname BookOriginal@termPointwiseNorm\endcsname}}
\expandafter\def\csname BookIndex@T0164\endcsname{suoboliefufanshu@索伯列夫范数（Sobolev norm）}
\expandafter\let\csname BookOriginal@termSobolevNorm\endcsname\termSobolevNorm
\renewcommand{\termSobolevNorm}{\BookTerm{T0164}{\csname BookOriginal@termSobolevNorm\endcsname}}
\expandafter\def\csname BookIndex@T0165\endcsname{jinhuayingshe@紧化映射（compactified map）}
\expandafter\let\csname BookOriginal@termCompactifiedMap\endcsname\termCompactifiedMap
\renewcommand{\termCompactifiedMap}{\BookTerm{T0165}{\csname BookOriginal@termCompactifiedMap\endcsname}}
\expandafter\def\csname BookIndex@T0166\endcsname{juanjiji@卷积积（convolution product）}
\expandafter\let\csname BookOriginal@termConvolutionProduct\endcsname\termConvolutionProduct
\renewcommand{\termConvolutionProduct}{\BookTerm{T0166}{\csname BookOriginal@termConvolutionProduct\endcsname}}
\expandafter\def\csname BookIndex@T0167\endcsname{yubianjieyingshe@预边界映射（pre boundary map）}
\expandafter\let\csname BookOriginal@termPreBoundaryMap\endcsname\termPreBoundaryMap
\renewcommand{\termPreBoundaryMap}{\BookTerm{T0167}{\csname BookOriginal@termPreBoundaryMap\endcsname}}
\expandafter\def\csname BookIndex@T0168\endcsname{yulianyingshe@预链映射（pre chain map）}
\expandafter\let\csname BookOriginal@termPreChainMap\endcsname\termPreChainMap
\renewcommand{\termPreChainMap}{\BookTerm{T0168}{\csname BookOriginal@termPreChainMap\endcsname}}
\expandafter\def\csname BookIndex@T0169\endcsname{zhishushuaijian@指数衰减（exponential decay）}
\expandafter\let\csname BookOriginal@termExponentialDecay\endcsname\termExponentialDecay
\renewcommand{\termExponentialDecay}{\BookTerm{T0169}{\csname BookOriginal@termExponentialDecay\endcsname}}
\expandafter\def\csname BookIndex@T0170\endcsname{xianyanguji@先验估计（apriori estimate）}
\expandafter\let\csname BookOriginal@termAprioriEstimate\endcsname\termAprioriEstimate
\renewcommand{\termAprioriEstimate}{\BookTerm{T0170}{\csname BookOriginal@termAprioriEstimate\endcsname}}
\expandafter\def\csname BookIndex@T0171\endcsname{shangbankongjian@上半空间（upper half space）}
\expandafter\let\csname BookOriginal@termUpperHalfSpace\endcsname\termUpperHalfSpace
\renewcommand{\termUpperHalfSpace}{\BookTerm{T0171}{\csname BookOriginal@termUpperHalfSpace\endcsname}}
\expandafter\def\csname BookIndex@T0172\endcsname{chazhidingli@插值定理（interpolation theorem）}
\expandafter\let\csname BookOriginal@termInterpolationTheorem\endcsname\termInterpolationTheorem
\renewcommand{\termInterpolationTheorem}{\BookTerm{T0172}{\csname BookOriginal@termInterpolationTheorem\endcsname}}
\expandafter\def\csname BookIndex@T0173\endcsname{cankaoduliang@参考度量（reference metric）}
\expandafter\let\csname BookOriginal@termReferenceMetric\endcsname\termReferenceMetric
\renewcommand{\termReferenceMetric}{\BookTerm{T0173}{\csname BookOriginal@termReferenceMetric\endcsname}}
\expandafter\def\csname BookIndex@T0174\endcsname{zilieshoulian@子列收敛（subsequence convergence）}
\expandafter\let\csname BookOriginal@termSubsequenceConvergence\endcsname\termSubsequenceConvergence
\renewcommand{\termSubsequenceConvergence}{\BookTerm{T0174}{\csname BookOriginal@termSubsequenceConvergence\endcsname}}
\expandafter\def\csname BookIndex@T0175\endcsname{jianjinsuanzi@渐近算子（asymptotic operator）}
\expandafter\let\csname BookOriginal@termAsymptoticOperator\endcsname\termAsymptoticOperator
\renewcommand{\termAsymptoticOperator}{\BookTerm{T0175}{\csname BookOriginal@termAsymptoticOperator\endcsname}}
\expandafter\def\csname BookIndex@T0176\endcsname{zibansuanzi@自伴算子（self adjoint operator）}
\expandafter\let\csname BookOriginal@termSelfAdjointOperator\endcsname\termSelfAdjointOperator
\renewcommand{\termSelfAdjointOperator}{\BookTerm{T0176}{\csname BookOriginal@termSelfAdjointOperator\endcsname}}
\expandafter\def\csname BookIndex@T0177\endcsname{lieweiqiweitalianluo@列维--奇维塔联络（levi civita connection）}
\expandafter\let\csname BookOriginal@termLeviCivitaConnection\endcsname\termLeviCivitaConnection
\renewcommand{\termLeviCivitaConnection}{\BookTerm{T0177}{\csname BookOriginal@termLeviCivitaConnection\endcsname}}
\expandafter\def\csname BookIndex@T0178\endcsname{quluzhangliang@曲率张量（curvature tensor）}
\expandafter\let\csname BookOriginal@termCurvatureTensor\endcsname\termCurvatureTensor
\renewcommand{\termCurvatureTensor}{\BookTerm{T0178}{\csname BookOriginal@termCurvatureTensor\endcsname}}
\expandafter\def\csname BookIndex@T0179\endcsname{zhengjiaotiaojian@正交条件（orthogonality condition）}
\expandafter\let\csname BookOriginal@termOrthogonalityCondition\endcsname\termOrthogonalityCondition
\renewcommand{\termOrthogonalityCondition}{\BookTerm{T0179}{\csname BookOriginal@termOrthogonalityCondition\endcsname}}
\expandafter\def\csname BookIndex@T0180\endcsname{l2fanshu@$L^2$-范数（ltwo norm）}
\expandafter\let\csname BookOriginal@termLtwoNorm\endcsname\termLtwoNorm
\renewcommand{\termLtwoNorm}{\BookTerm{T0180}{\csname BookOriginal@termLtwoNorm\endcsname}}
\expandafter\def\csname BookIndex@T0181\endcsname{l2neiji@$L^2$-内积（ltwo inner product）}
\expandafter\let\csname BookOriginal@termLtwoInnerProduct\endcsname\termLtwoInnerProduct
\renewcommand{\termLtwoInnerProduct}{\BookTerm{T0181}{\csname BookOriginal@termLtwoInnerProduct\endcsname}}
\expandafter\def\csname BookIndex@T0182\endcsname{pingtanzhumianduliang@平坦柱面度量（flat cylindrical metric）}
\expandafter\let\csname BookOriginal@termFlatCylindricalMetric\endcsname\termFlatCylindricalMetric
\renewcommand{\termFlatCylindricalMetric}{\BookTerm{T0182}{\csname BookOriginal@termFlatCylindricalMetric\endcsname}}
\expandafter\def\csname BookIndex@T0183\endcsname{dengwenzuobiao@等温坐标（isothermal coordinates）}
\expandafter\let\csname BookOriginal@termIsothermalCoordinates\endcsname\termIsothermalCoordinates
\renewcommand{\termIsothermalCoordinates}{\BookTerm{T0183}{\csname BookOriginal@termIsothermalCoordinates\endcsname}}
\expandafter\def\csname BookIndex@T0184\endcsname{erjiejibenxingshi@二阶基本形式（second fundamental form）}
\expandafter\let\csname BookOriginal@termSecondFundamentalForm\endcsname\termSecondFundamentalForm
\renewcommand{\termSecondFundamentalForm}{\BookTerm{T0184}{\csname BookOriginal@termSecondFundamentalForm\endcsname}}
\expandafter\def\csname BookIndex@T0185\endcsname{quancede@全测地（totally geodesic）}
\expandafter\let\csname BookOriginal@termTotallyGeodesic\endcsname\termTotallyGeodesic
\renewcommand{\termTotallyGeodesic}{\BookTerm{T0185}{\csname BookOriginal@termTotallyGeodesic\endcsname}}
\expandafter\def\csname BookIndex@T0186\endcsname{bixiang@闭像（closed range）}
\expandafter\let\csname BookOriginal@termClosedRange\endcsname\termClosedRange
\renewcommand{\termClosedRange}{\BookTerm{T0186}{\csname BookOriginal@termClosedRange\endcsname}}
\expandafter\def\csname BookIndex@T0187\endcsname{youpingfanhua@酉平凡化（unitary trivialization）}
\expandafter\let\csname BookOriginal@termUnitaryTrivialization\endcsname\termUnitaryTrivialization
\renewcommand{\termUnitaryTrivialization}{\BookTerm{T0187}{\csname BookOriginal@termUnitaryTrivialization\endcsname}}
\expandafter\def\csname BookIndex@T0188\endcsname{jinsuanzi@紧算子（compact operator）}
\expandafter\let\csname BookOriginal@termCompactOperator\endcsname\termCompactOperator
\renewcommand{\termCompactOperator}{\BookTerm{T0188}{\csname BookOriginal@termCompactOperator\endcsname}}
\expandafter\def\csname BookIndex@T0189\endcsname{kaiyingshedingli@开映射定理（open mapping theorem）}
\expandafter\let\csname BookOriginal@termOpenMappingTheorem\endcsname\termOpenMappingTheorem
\renewcommand{\termOpenMappingTheorem}{\BookTerm{T0189}{\csname BookOriginal@termOpenMappingTheorem\endcsname}}
\expandafter\def\csname BookIndex@T0190\endcsname{zuidazhiyuanli@最大值原理（maximum principle）}
\expandafter\let\csname BookOriginal@termMaximumPrinciple\endcsname\termMaximumPrinciple
\renewcommand{\termMaximumPrinciple}{\BookTerm{T0190}{\csname BookOriginal@termMaximumPrinciple\endcsname}}
\expandafter\def\csname BookIndex@T0191\endcsname{suoboliefuqianru@索伯列夫嵌入（Sobolev embedding）}
\expandafter\let\csname BookOriginal@termSobolevEmbedding\endcsname\termSobolevEmbedding
\renewcommand{\termSobolevEmbedding}{\BookTerm{T0191}{\csname BookOriginal@termSobolevEmbedding\endcsname}}
\expandafter\def\csname BookIndex@T0192\endcsname{kaileijiao@凯勒角（Kähler angle）}
\expandafter\let\csname BookOriginal@termKahlerAngle\endcsname\termKahlerAngle
\renewcommand{\termKahlerAngle}{\BookTerm{T0192}{\csname BookOriginal@termKahlerAngle\endcsname}}
\expandafter\def\csname BookIndex@T0193\endcsname{fenlieduan@分裂端（splitting end）}
\expandafter\let\csname BookOriginal@termSplittingEnd\endcsname\termSplittingEnd
\renewcommand{\termSplittingEnd}{\BookTerm{T0193}{\csname BookOriginal@termSplittingEnd\endcsname}}
\expandafter\def\csname BookIndex@T0194\endcsname{maopaoduan@冒泡端（bubbling end）}
\expandafter\let\csname BookOriginal@termBubblingEnd\endcsname\termBubblingEnd
\renewcommand{\termBubblingEnd}{\BookTerm{T0194}{\csname BookOriginal@termBubblingEnd\endcsname}}
\expandafter\def\csname BookIndex@T0195\endcsname{shangtapu@商拓扑（quotient topology）}
\expandafter\let\csname BookOriginal@termQuotientTopology\endcsname\termQuotientTopology
\renewcommand{\termQuotientTopology}{\BookTerm{T0195}{\csname BookOriginal@termQuotientTopology\endcsname}}
\expandafter\def\csname BookIndex@T0196\endcsname{jingxicinftytapu@精细 $C^\infty$ 拓扑（fine c\ensuremath{\infty} topology）}
\expandafter\let\csname BookOriginal@termFineCinfinityTopology\endcsname\termFineCinfinityTopology
\renewcommand{\termFineCinfinityTopology}{\BookTerm{T0196}{\csname BookOriginal@termFineCinfinityTopology\endcsname}}
\expandafter\def\csname BookIndex@T0197\endcsname{duanliejiandianliuxian@断裂尖点流线（broken cusp trajectory）}
\expandafter\let\csname BookOriginal@termBrokenCuspTrajectory\endcsname\termBrokenCuspTrajectory
\renewcommand{\termBrokenCuspTrajectory}{\BookTerm{T0197}{\csname BookOriginal@termBrokenCuspTrajectory\endcsname}}
\expandafter\def\csname BookIndex@T0198\endcsname{feidingchang@非定常（non stationary）}
\expandafter\let\csname BookOriginal@termNonStationary\endcsname\termNonStationary
\renewcommand{\termNonStationary}{\BookTerm{T0198}{\csname BookOriginal@termNonStationary\endcsname}}
\expandafter\def\csname BookIndex@T0199\endcsname{huifuxulie@恢复序列（recovering sequence）}
\expandafter\let\csname BookOriginal@termRecoveringSequence\endcsname\termRecoveringSequence
\renewcommand{\termRecoveringSequence}{\BookTerm{T0199}{\csname BookOriginal@termRecoveringSequence\endcsname}}
\expandafter\def\csname BookIndex@T0200\endcsname{jidalian@极大链（maximal chain）}
\expandafter\let\csname BookOriginal@termMaximalChain\endcsname\termMaximalChain
\renewcommand{\termMaximalChain}{\BookTerm{T0200}{\csname BookOriginal@termMaximalChain\endcsname}}
\expandafter\def\csname BookIndex@T0201\endcsname{nengliangguiyihua@能量归一化（energy normalization）}
\expandafter\let\csname BookOriginal@termEnergyNormalization\endcsname\termEnergyNormalization
\renewcommand{\termEnergyNormalization}{\BookTerm{T0201}{\csname BookOriginal@termEnergyNormalization\endcsname}}
\expandafter\def\csname BookIndex@T0202\endcsname{qiepian@切片（slice）}
\expandafter\let\csname BookOriginal@termSlice\endcsname\termSlice
\renewcommand{\termSlice}{\BookTerm{T0202}{\csname BookOriginal@termSlice\endcsname}}
\expandafter\def\csname BookIndex@T0203\endcsname{linyuji@邻域基（neighborhood basis）}
\expandafter\let\csname BookOriginal@termNeighborhoodBasis\endcsname\termNeighborhoodBasis
\renewcommand{\termNeighborhoodBasis}{\BookTerm{T0203}{\csname BookOriginal@termNeighborhoodBasis\endcsname}}
\expandafter\def\csname BookIndex@T0204\endcsname{epsilonzhengzexingdingli@$\epsilon$-正则性定理（epsilon regularity theorem）}
\expandafter\let\csname BookOriginal@termEpsilonRegularityTheorem\endcsname\termEpsilonRegularityTheorem
\renewcommand{\termEpsilonRegularityTheorem}{\BookTerm{T0204}{\csname BookOriginal@termEpsilonRegularityTheorem\endcsname}}
\expandafter\def\csname BookIndex@T0205\endcsname{sakesiwulunbeikemaopaolunzheng@萨克斯--乌伦贝克冒泡论证（sacks uhlenbeck bubbling argument）}
\expandafter\let\csname BookOriginal@termSacksUhlenbeckBubblingArgument\endcsname\termSacksUhlenbeckBubblingArgument
\renewcommand{\termSacksUhlenbeckBubblingArgument}{\BookTerm{T0205}{\csname BookOriginal@termSacksUhlenbeckBubblingArgument\endcsname}}
\expandafter\def\csname BookIndex@T0206\endcsname{asikeliaerzeladingli@阿斯科利--阿尔泽拉定理（Ascoli Arzelà theorem）}
\expandafter\let\csname BookOriginal@termAscoliArzelaTheorem\endcsname\termAscoliArzelaTheorem
\renewcommand{\termAscoliArzelaTheorem}{\BookTerm{T0206}{\csname BookOriginal@termAscoliArzelaTheorem\endcsname}}
\expandafter\def\csname BookIndex@T0207\endcsname{leibeigekongzhishouliandingli@勒贝格控制收敛定理（lebesgue dominated convergence theorem）}
\expandafter\let\csname BookOriginal@termLebesgueDominatedConvergenceTheorem\endcsname\termLebesgueDominatedConvergenceTheorem
\renewcommand{\termLebesgueDominatedConvergenceTheorem}{\BookTerm{T0207}{\csname BookOriginal@termLebesgueDominatedConvergenceTheorem\endcsname}}
\expandafter\def\csname BookIndex@T0208\endcsname{zhufenzhi@主分支（principal component）}
\expandafter\let\csname BookOriginal@termPrincipalComponent\endcsname\termPrincipalComponent
\renewcommand{\termPrincipalComponent}{\BookTerm{T0208}{\csname BookOriginal@termPrincipalComponent\endcsname}}
\expandafter\def\csname BookIndex@T0209\endcsname{qipaofenzhi@气泡分支（bubble component）}
\expandafter\let\csname BookOriginal@termBubbleComponent\endcsname\termBubbleComponent
\renewcommand{\termBubbleComponent}{\BookTerm{T0209}{\csname BookOriginal@termBubbleComponent\endcsname}}
\expandafter\def\csname BookIndex@T0210\endcsname{bukeyuefenzhi@不可约分支（irreducible component）}
\expandafter\let\csname BookOriginal@termIrreducibleComponent\endcsname\termIrreducibleComponent
\renewcommand{\termIrreducibleComponent}{\BookTerm{T0210}{\csname BookOriginal@termIrreducibleComponent\endcsname}}
\expandafter\def\csname BookIndex@T0211\endcsname{yaobanqunzuoyong@幺半群作用（monoidal action）}
\expandafter\let\csname BookOriginal@termMonoidalAction\endcsname\termMonoidalAction
\renewcommand{\termMonoidalAction}{\BookTerm{T0211}{\csname BookOriginal@termMonoidalAction\endcsname}}
\expandafter\def\csname BookIndex@T0212\endcsname{shuangyaobanqunzuoyong@双幺半群作用（bi monoidal action）}
\expandafter\let\csname BookOriginal@termBiMonoidalAction\endcsname\termBiMonoidalAction
\renewcommand{\termBiMonoidalAction}{\BookTerm{T0212}{\csname BookOriginal@termBiMonoidalAction\endcsname}}
\expandafter\def\csname BookIndex@T0213\endcsname{wendingyingshetapu@稳定映射拓扑（stable map topology）}
\expandafter\let\csname BookOriginal@termStableMapTopology\endcsname\termStableMapTopology
\renewcommand{\termStableMapTopology}{\BookTerm{T0213}{\csname BookOriginal@termStableMapTopology\endcsname}}
\expandafter\def\csname BookIndex@T0214\endcsname{daibianwendingyingshe@带边稳定映射（bordered stable map）}
\expandafter\let\csname BookOriginal@termBorderedStableMap\endcsname\termBorderedStableMap
\renewcommand{\termBorderedStableMap}{\BookTerm{T0214}{\csname BookOriginal@termBorderedStableMap\endcsname}}
\expandafter\def\csname BookIndex@T0215\endcsname{xianweiji@纤维积（fiber product）}
\expandafter\let\csname BookOriginal@termFiberProduct\endcsname\termFiberProduct
\renewcommand{\termFiberProduct}{\BookTerm{T0215}{\csname BookOriginal@termFiberProduct\endcsname}}
\expandafter\def\csname BookIndex@T0216\endcsname{pipeitiaojian@匹配条件（matching condition）}
\expandafter\let\csname BookOriginal@termMatchingCondition\endcsname\termMatchingCondition
\renewcommand{\termMatchingCondition}{\BookTerm{T0216}{\csname BookOriginal@termMatchingCondition\endcsname}}
\expandafter\def\csname BookIndex@T0217\endcsname{banwendingquxian@半稳定曲线（semi stable curve）}
\expandafter\let\csname BookOriginal@termSemiStableCurve\endcsname\termSemiStableCurve
\renewcommand{\termSemiStableCurve}{\BookTerm{T0217}{\csname BookOriginal@termSemiStableCurve\endcsname}}
\expandafter\def\csname BookIndex@T0218\endcsname{qipaogendian@气泡根点（bubble root）}
\expandafter\let\csname BookOriginal@termBubbleRoot\endcsname\termBubbleRoot
\renewcommand{\termBubbleRoot}{\BookTerm{T0218}{\csname BookOriginal@termBubbleRoot\endcsname}}
\expandafter\def\csname BookIndex@T0219\endcsname{fenbuyingshe@分布映射（distribution map）}
\expandafter\let\csname BookOriginal@termDistributionMap\endcsname\termDistributionMap
\renewcommand{\termDistributionMap}{\BookTerm{T0219}{\csname BookOriginal@termDistributionMap\endcsname}}
\expandafter\def\csname BookIndex@T0220\endcsname{yuji@余积（coproduct）}
\expandafter\let\csname BookOriginal@termCoproduct\endcsname\termCoproduct
\renewcommand{\termCoproduct}{\BookTerm{T0220}{\csname BookOriginal@termCoproduct\endcsname}}
\expandafter\def\csname BookIndex@T0221\endcsname{qiangcinftytapu@强 $C^\infty$ 拓扑（strong c\ensuremath{\infty} topology）}
\expandafter\let\csname BookOriginal@termStrongCinfinityTopology\endcsname\termStrongCinfinityTopology
\renewcommand{\termStrongCinfinityTopology}{\BookTerm{T0221}{\csname BookOriginal@termStrongCinfinityTopology\endcsname}}
\expandafter\def\csname BookIndex@T0222\endcsname{daijiaoorbifold@带角 orbifold（orbifold with corners）}
\expandafter\let\csname BookOriginal@termOrbifoldWithCorners\endcsname\termOrbifoldWithCorners
\renewcommand{\termOrbifoldWithCorners}{\BookTerm{T0222}{\csname BookOriginal@termOrbifoldWithCorners\endcsname}}
\expandafter\def\csname BookIndex@T0223\endcsname{lianyingshemokongjian@链映射模空间（chain map moduli space）}
\expandafter\let\csname BookOriginal@termChainMapModuliSpace\endcsname\termChainMapModuliSpace
\renewcommand{\termChainMapModuliSpace}{\BookTerm{T0223}{\csname BookOriginal@termChainMapModuliSpace\endcsname}}
\expandafter\def\csname BookIndex@T0224\endcsname{nengliangguji@能量估计（energy estimate）}
\expandafter\let\csname BookOriginal@termEnergyEstimate\endcsname\termEnergyEstimate
\renewcommand{\termEnergyEstimate}{\BookTerm{T0224}{\csname BookOriginal@termEnergyEstimate\endcsname}}
\expandafter\def\csname BookIndex@T0225\endcsname{luguobianhua@滤过变化（filtration change）}
\expandafter\let\csname BookOriginal@termFiltrationChange\endcsname\termFiltrationChange
\renewcommand{\termFiltrationChange}{\BookTerm{T0225}{\csname BookOriginal@termFiltrationChange\endcsname}}
\expandafter\def\csname BookIndex@T0226\endcsname{zuobiaobianhuan@坐标变换（coordinate change）}
\expandafter\let\csname BookOriginal@termCoordinateChange\endcsname\termCoordinateChange
\renewcommand{\termCoordinateChange}{\BookTerm{T0226}{\csname BookOriginal@termCoordinateChange\endcsname}}
\expandafter\def\csname BookIndex@T0227\endcsname{hamidunliangdexingbian@哈密顿量的形变（deformation of Hamiltonians）}
\expandafter\let\csname BookOriginal@termDeformationOfHamiltonians\endcsname\termDeformationOfHamiltonians
\renewcommand{\termDeformationOfHamiltonians}{\BookTerm{T0227}{\csname BookOriginal@termDeformationOfHamiltonians\endcsname}}
\expandafter\def\csname BookIndex@T0228\endcsname{fuluoeryantayingshe@弗洛尔延拓映射（Floer continuation map）}
\expandafter\let\csname BookOriginal@termFloerContinuationMap\endcsname\termFloerContinuationMap
\renewcommand{\termFloerContinuationMap}{\BookTerm{T0228}{\csname BookOriginal@termFloerContinuationMap\endcsname}}
\expandafter\def\csname BookIndex@T0229\endcsname{cengjibianhua@层级变化（level change）}
\expandafter\let\csname BookOriginal@termLevelChange\endcsname\termLevelChange
\renewcommand{\termLevelChange}{\BookTerm{T0229}{\csname BookOriginal@termLevelChange\endcsname}}
\expandafter\def\csname BookIndex@T0230\endcsname{tongshihamiduntonghen@同时哈密顿同痕（simultaneous Hamiltonian isotopy）}
\expandafter\let\csname BookOriginal@termSimultaneousHamiltonianIsotopy\endcsname\termSimultaneousHamiltonianIsotopy
\renewcommand{\termSimultaneousHamiltonianIsotopy}{\BookTerm{T0230}{\csname BookOriginal@termSimultaneousHamiltonianIsotopy\endcsname}}
\expandafter\def\csname BookIndex@T0231\endcsname{jianjintiaojian@渐近条件（asymptotic condition）}
\expandafter\let\csname BookOriginal@termAsymptoticCondition\endcsname\termAsymptoticCondition
\renewcommand{\termAsymptoticCondition}{\BookTerm{T0231}{\csname BookOriginal@termAsymptoticCondition\endcsname}}
\expandafter\def\csname BookIndex@T0232\endcsname{tongluntiaojian@同伦条件（homotopy condition）}
\expandafter\let\csname BookOriginal@termHomotopyCondition\endcsname\termHomotopyCondition
\renewcommand{\termHomotopyCondition}{\BookTerm{T0232}{\csname BookOriginal@termHomotopyCondition\endcsname}}
\expandafter\def\csname BookIndex@T0233\endcsname{xingbianfuxing@形变复形（deformation complex）}
\expandafter\let\csname BookOriginal@termDeformationComplex\endcsname\termDeformationComplex
\renewcommand{\termDeformationComplex}{\BookTerm{T0233}{\csname BookOriginal@termDeformationComplex\endcsname}}
\expandafter\def\csname BookIndex@T0234\endcsname{banaheliuxing@巴拿赫流形（Banach manifold）}
\expandafter\let\csname BookOriginal@termBanachManifold\endcsname\termBanachManifold
\renewcommand{\termBanachManifold}{\BookTerm{T0234}{\csname BookOriginal@termBanachManifold\endcsname}}
\expandafter\def\csname BookIndex@T0235\endcsname{jiaquansuoboliefukongjian@加权索伯列夫空间（weighted Sobolev space）}
\expandafter\let\csname BookOriginal@termWeightedSobolevSpace\endcsname\termWeightedSobolevSpace
\renewcommand{\termWeightedSobolevSpace}{\BookTerm{T0235}{\csname BookOriginal@termWeightedSobolevSpace\endcsname}}
\expandafter\def\csname BookIndex@T0236\endcsname{jiandianliuxiankongjian@尖点流线空间（cusp trajectory space）}
\expandafter\let\csname BookOriginal@termCuspTrajectorySpace\endcsname\termCuspTrajectorySpace
\renewcommand{\termCuspTrajectorySpace}{\BookTerm{T0236}{\csname BookOriginal@termCuspTrajectorySpace\endcsname}}
\expandafter\def\csname BookIndex@T0237\endcsname{zuzhibiaolilun@族指标理论（family index theory）}
\expandafter\let\csname BookOriginal@termFamilyIndexTheory\endcsname\termFamilyIndexTheory
\renewcommand{\termFamilyIndexTheory}{\BookTerm{T0237}{\csname BookOriginal@termFamilyIndexTheory\endcsname}}
\expandafter\def\csname BookIndex@T0238\endcsname{xinglieshixiancong@行列式线丛（determinant line bundle）}
\expandafter\let\csname BookOriginal@termDeterminantLineBundle\endcsname\termDeterminantLineBundle
\renewcommand{\termDeterminantLineBundle}{\BookTerm{T0238}{\csname BookOriginal@termDeterminantLineBundle\endcsname}}
\expandafter\def\csname BookIndex@T0239\endcsname{xiangduispinjiegou@相对 spin 结构（relative spin structure）}
\expandafter\let\csname BookOriginal@termRelativeSpinStructure\endcsname\termRelativeSpinStructure
\renewcommand{\termRelativeSpinStructure}{\BookTerm{T0239}{\csname BookOriginal@termRelativeSpinStructure\endcsname}}
\expandafter\def\csname BookIndex@T0240\endcsname{xiangduispin@相对 spin（relatively spin）}
\expandafter\let\csname BookOriginal@termRelativelySpin\endcsname\termRelativelySpin
\renewcommand{\termRelativelySpin}{\BookTerm{T0240}{\csname BookOriginal@termRelativelySpin\endcsname}}
\expandafter\def\csname BookIndex@T0241\endcsname{jinsijie@近似解（approximate solution）}
\expandafter\let\csname BookOriginal@termApproximateSolution\endcsname\termApproximateSolution
\renewcommand{\termApproximateSolution}{\BookTerm{T0241}{\csname BookOriginal@termApproximateSolution\endcsname}}
\expandafter\def\csname BookIndex@T0242\endcsname{youni@右逆（right inverse）}
\expandafter\let\csname BookOriginal@termRightInverse\endcsname\termRightInverse
\renewcommand{\termRightInverse}{\BookTerm{T0242}{\csname BookOriginal@termRightInverse\endcsname}}
\expandafter\def\csname BookIndex@T0243\endcsname{yasuoyingshedingli@压缩映射定理（contraction mapping theorem）}
\expandafter\let\csname BookOriginal@termContractionMappingTheorem\endcsname\termContractionMappingTheorem
\renewcommand{\termContractionMappingTheorem}{\BookTerm{T0243}{\csname BookOriginal@termContractionMappingTheorem\endcsname}}
\expandafter\def\csname BookIndex@T0244\endcsname{yuzhanheyingshe@预粘合映射（pre gluing map）}
\expandafter\let\csname BookOriginal@termPreGluingMap\endcsname\termPreGluingMap
\renewcommand{\termPreGluingMap}{\BookTerm{T0244}{\csname BookOriginal@termPreGluingMap\endcsname}}
\expandafter\def\csname BookIndex@T0245\endcsname{zhanheyingshe@粘合映射（gluing map）}
\expandafter\let\csname BookOriginal@termGluingMap\endcsname\termGluingMap
\renewcommand{\termGluingMap}{\BookTerm{T0245}{\csname BookOriginal@termGluingMap\endcsname}}
\expandafter\def\csname BookIndex@T0246\endcsname{lingkoulinyu@领口邻域（collar neighborhood）}
\expandafter\let\csname BookOriginal@termCollarNeighborhood\endcsname\termCollarNeighborhood
\renewcommand{\termCollarNeighborhood}{\BookTerm{T0246}{\csname BookOriginal@termCollarNeighborhood\endcsname}}
\expandafter\def\csname BookIndex@T0247\endcsname{fanchang@反常（anomaly）}
\expandafter\let\csname BookOriginal@termAnomaly\endcsname\termAnomaly
\renewcommand{\termAnomaly}{\BookTerm{T0247}{\csname BookOriginal@termAnomaly\endcsname}}
\expandafter\def\csname BookIndex@T0248\endcsname{shunzixiaoying@瞬子效应（instanton effects）}
\expandafter\let\csname BookOriginal@termInstantonEffects\endcsname\termInstantonEffects
\renewcommand{\termInstantonEffects}{\BookTerm{T0248}{\csname BookOriginal@termInstantonEffects\endcsname}}
\expandafter\def\csname BookIndex@T0249\endcsname{duozhijiemian@多值截面（multi valued section）}
\expandafter\let\csname BookOriginal@termMultiValuedSection\endcsname\termMultiValuedSection
\renewcommand{\termMultiValuedSection}{\BookTerm{T0249}{\csname BookOriginal@termMultiValuedSection\endcsname}}
\expandafter\def\csname BookIndex@T0250\endcsname{hengjiedui@横截对（transverse pair）}
\expandafter\let\csname BookOriginal@termTransversePair\endcsname\termTransversePair
\renewcommand{\termTransversePair}{\BookTerm{T0250}{\csname BookOriginal@termTransversePair\endcsname}}
\expandafter\def\csname BookIndex@T0251\endcsname{dandiaochangshu@单调常数（monotonicity constant）}
\expandafter\let\csname BookOriginal@termMonotonicityConstant\endcsname\termMonotonicityConstant
\renewcommand{\termMonotonicityConstant}{\BookTerm{T0251}{\csname BookOriginal@termMonotonicityConstant\endcsname}}
\expandafter\def\csname BookIndex@T0252\endcsname{zaimouchudanshe@在某处单射（somewhere injective）}
\expandafter\let\csname BookOriginal@termSomewhereInjective\endcsname\termSomewhereInjective
\renewcommand{\termSomewhereInjective}{\BookTerm{T0252}{\csname BookOriginal@termSomewhereInjective\endcsname}}
\expandafter\def\csname BookIndex@T0253\endcsname{sanjiaoxingchengji@三角形乘积（triangle product）}
\expandafter\let\csname BookOriginal@termTriangleProduct\endcsname\termTriangleProduct
\renewcommand{\termTriangleProduct}{\BookTerm{T0253}{\csname BookOriginal@termTriangleProduct\endcsname}}
\expandafter\def\csname BookIndex@T0254\endcsname{luguosanjiaoxingchengji@滤过三角形乘积（filtered triangle product）}
\expandafter\let\csname BookOriginal@termFilteredTriangleProduct\endcsname\termFilteredTriangleProduct
\renewcommand{\termFilteredTriangleProduct}{\BookTerm{T0254}{\csname BookOriginal@termFilteredTriangleProduct\endcsname}}
\expandafter\def\csname BookIndex@T0255\endcsname{yilaidingyiyudiandezu@依赖定义域点的族（domain dependent family）}
\expandafter\let\csname BookOriginal@termDomainDependentFamily\endcsname\termDomainDependentFamily
\renewcommand{\termDomainDependentFamily}{\BookTerm{T0255}{\csname BookOriginal@termDomainDependentFamily\endcsname}}
\expandafter\def\csname BookIndex@T0256\endcsname{vtunmeihamiduntonghen@$V$-吞没哈密顿同痕（v engulfed Hamiltonian isotopy）}
\expandafter\let\csname BookOriginal@termVEngulfedHamiltonianIsotopy\endcsname\termVEngulfedHamiltonianIsotopy
\renewcommand{\termVEngulfedHamiltonianIsotopy}{\BookTerm{T0256}{\csname BookOriginal@termVEngulfedHamiltonianIsotopy\endcsname}}
\expandafter\def\csname BookIndex@T0257\endcsname{houbaoerfen@厚--薄二分（thick thin dichotomy）}
\expandafter\let\csname BookOriginal@termThickThinDichotomy\endcsname\termThickThinDichotomy
\renewcommand{\termThickThinDichotomy}{\BookTerm{T0257}{\csname BookOriginal@termThickThinDichotomy\endcsname}}
\expandafter\def\csname BookIndex@T0258\endcsname{wuqiongyuanchutu@无穷远处凸（convex at \ensuremath{\infty}）}
\expandafter\let\csname BookOriginal@termConvexAtInfinity\endcsname\termConvexAtInfinity
\renewcommand{\termConvexAtInfinity}{\BookTerm{T0258}{\csname BookOriginal@termConvexAtInfinity\endcsname}}
\expandafter\def\csname BookIndex@T0259\endcsname{qiangxintianchong@强辛填充（strong symplectic filling）}
\expandafter\let\csname BookOriginal@termStrongSymplecticFilling\endcsname\termStrongSymplecticFilling
\renewcommand{\termStrongSymplecticFilling}{\BookTerm{T0259}{\csname BookOriginal@termStrongSymplecticFilling\endcsname}}
\expandafter\def\csname BookIndex@T0260\endcsname{zhumianzuobiao@柱面坐标（cylindrical coordinates）}
\expandafter\let\csname BookOriginal@termCylindricalCoordinates\endcsname\termCylindricalCoordinates
\renewcommand{\termCylindricalCoordinates}{\BookTerm{T0260}{\csname BookOriginal@termCylindricalCoordinates\endcsname}}
\expandafter\def\csname BookIndex@T0261\endcsname{zhuizuobiao@锥坐标（cone coordinates）}
\expandafter\let\csname BookOriginal@termConeCoordinates\endcsname\termConeCoordinates
\renewcommand{\termConeCoordinates}{\BookTerm{T0261}{\csname BookOriginal@termConeCoordinates\endcsname}}
\expandafter\def\csname BookIndex@T0262\endcsname{jiechusanyuanzu@接触三元组（contact triad）}
\expandafter\let\csname BookOriginal@termContactTriad\endcsname\termContactTriad
\renewcommand{\termContactTriad}{\BookTerm{T0262}{\csname BookOriginal@termContactTriad\endcsname}}
\expandafter\def\csname BookIndex@T0263\endcsname{sanyuanzuduliang@三元组度量（triad metric）}
\expandafter\let\csname BookOriginal@termTriadMetric\endcsname\termTriadMetric
\renewcommand{\termTriadMetric}{\BookTerm{T0263}{\csname BookOriginal@termTriadMetric\endcsname}}
\expandafter\def\csname BookIndex@T0264\endcsname{lambdajiechuxing@$\lambda$-接触型（lambda contact type）}
\expandafter\let\csname BookOriginal@termLambdaContactType\endcsname\termLambdaContactType
\renewcommand{\termLambdaContactType}{\BookTerm{T0264}{\csname BookOriginal@termLambdaContactType\endcsname}}
\expandafter\def\csname BookIndex@T0265\endcsname{cilinjie@次临界（subcritical）}
\expandafter\let\csname BookOriginal@termSubcritical\endcsname\termSubcritical
\renewcommand{\termSubcritical}{\BookTerm{T0265}{\csname BookOriginal@termSubcritical\endcsname}}
\expandafter\def\csname BookIndex@T0266\endcsname{shangtongdiaoqiumianxing@上同调球面型（cohomologically spherical）}
\expandafter\let\csname BookOriginal@termCohomologicallySpherical\endcsname\termCohomologicallySpherical
\renewcommand{\termCohomologicallySpherical}{\BookTerm{T0266}{\csname BookOriginal@termCohomologicallySpherical\endcsname}}
\expandafter\def\csname BookIndex@T0267\endcsname{wendinglagelangrijiaoxingzhi@稳定拉格朗日交性质（stable Lagrangian intersection property）}
\expandafter\let\csname BookOriginal@termStableLagrangianIntersectionProperty\endcsname\termStableLagrangianIntersectionProperty
\renewcommand{\termStableLagrangianIntersectionProperty}{\BookTerm{T0267}{\csname BookOriginal@termStableLagrangianIntersectionProperty\endcsname}}
\expandafter\def\csname BookIndex@T0268\endcsname{tianchongyuanpan@填充圆盘（bounding disc）}
\expandafter\let\csname BookOriginal@termBoundingDisc\endcsname\termBoundingDisc
\renewcommand{\termBoundingDisc}{\BookTerm{T0268}{\csname BookOriginal@termBoundingDisc\endcsname}}
\expandafter\def\csname BookIndex@T0269\endcsname{dingchangjie@定常解（stationary solution）}
\expandafter\let\csname BookOriginal@termStationarySolution\endcsname\termStationarySolution
\renewcommand{\termStationarySolution}{\BookTerm{T0269}{\csname BookOriginal@termStationarySolution\endcsname}}
\makeatother
